\section*{Разбор задачи из тестового варианта №2}
\addcontentsline{toc}{section}{Разбор задачи из тестового варианта №2}
\begin{Task}{Сходимость по распределению и Дельта-метод}{task2}
Случайные величины $(X_n, n \in \mathbb{N})$ независимы и имеют экспоненциальное распределение $\text{Exp}(\theta)$, $\theta > 0$. Обозначим $\overline{X} = \frac{1}{n} \sum_{i=1}^n X_i$. Найдите предел сходимости по распределению у последовательности
\[
\sqrt{n} \left( \frac{1}{\overline{X}^2} - \theta^2 \right).
\]
\end{Task}

\begin{Solution}
Задача требует найти асимптотическое распределение гладкой нелинейной функции от выборочного среднего. Прямое вычисление распределения $\overline{X}$ (которое имеет гамма-распределение) с последующим преобразованием приведет к неоправданно сложным вычислениям. Адекватным математическим аппаратом для данного класса задач является комбинация Центральной предельной теоремы (ЦПТ) и теоремы о наследовании сходимости (Дельта-метода).

\begin{MethodBox}[Одномерный Дельта-метод]
Пусть последовательность случайных величин $Y_n$ (чаще всего $Y_n = \overline{X}$) удовлетворяет условию асимптотической нормальности:
\[
\sqrt{n}(Y_n - \mu) \toD \N(0, \sigma^2).
\]
Тогда для любой функции $g(t)$, дифференцируемой в точке $\mu$, причем $g'(\mu) \neq 0$, справедливо соотношение:
\[
\sqrt{n}(g(Y_n) - g(\mu)) \toD \N(0, \sigma_{as}^2),
\]
где асимптотическая дисперсия вычисляется как $\sigma_{as}^2 = [g'(\mu)]^2 \cdot \sigma^2$.
\end{MethodBox}

\textbf{Шаг 1. Характеристики базового распределения.}\\
Рассмотрим элементы исходной выборки $X_i \sim \text{Exp}(\theta)$. Математическое ожидание и дисперсия экспоненциального распределения равны:
\[
\mu = \E X_1 = \frac{1}{\theta}, \quad \sigma^2 = \Var X_1 = \frac{1}{\theta^2}.
\]

\textbf{Шаг 2. Применение Центральной предельной теоремы.}\\
Поскольку случайные величины $X_i$ независимы и одинаково распределены, а их дисперсия конечна, по классической ЦПТ последовательность выборочных средних $\overline{X}$ является асимптотически нормальной:
\[
\sqrt{n} \left( \overline{X} - \frac{1}{\theta} \right) \toD \N\left(0, \frac{1}{\theta^2}\right).
\]

\textbf{Шаг 3. Анализ функции преобразования.}\\
Исследуемая последовательность имеет вид $\sqrt{n}(g(\overline{X}) - C)$. Введем функцию:
\[
g(t) = \frac{1}{t^2} = t^{-2}.
\]
Проверим корректность константы $C$ в условии задачи. Согласно Дельта-методу, вычитаемая константа должна строго равняться значению функции в точке математического ожидания:
\[
g(\mu) = g\left(\frac{1}{\theta}\right) = \frac{1}{(1/\theta)^2} = \theta^2.
\]
Это в точности совпадает с выражением $\theta^2$ из условия, что подтверждает применимость метода.

\textbf{Шаг 4. Вычисление асимптотической дисперсии.}\\
Найдем производную функции $g(t)$ и вычислим ее значение в точке $t = \mu = 1/\theta$:
\[
g'(t) = -2t^{-3} \implies g'\left(\frac{1}{\theta}\right) = -2\left(\frac{1}{\theta}\right)^{-3} = -2\theta^3.
\]
Заметим, что $g'(1/\theta) \neq 0$ при $\theta > 0$, следовательно, классический Дельта-метод (первого порядка) применим. Вычислим итоговую асимптотическую дисперсию:
\[
\sigma_{as}^2 = [g'(\mu)]^2 \cdot \Var X_1 = (-2\theta^3)^2 \cdot \frac{1}{\theta^2} = 4\theta^6 \cdot \frac{1}{\theta^2} = 4\theta^4.
\]

\begin{WarningBox}[Параметризация распределения и типичные ошибки]
Критически важно не путать параметризацию экспоненциального распределения. Если бы в условии было задано $X \sim \text{Exp}(\lambda)$ с плотностью $p(x) = \frac{1}{\lambda} e^{-x/\lambda}$, то математическое ожидание равнялось бы $\lambda$. В стандартной параметризации через интенсивность (как в данной задаче) $\E X = 1/\theta$. Также частой ошибкой является потеря квадрата при возведении производной: дисперсия масштабируется с квадратичной скоростью, а не линейной.
\end{WarningBox}

Синтезируя полученные результаты, приходим к выводу, что исходная последовательность сходится по распределению к центрированной нормальной величине с дисперсией $4\theta^4$.
\end{Solution}

\begin{Answer}
$\N(0, 4\theta^4)$
\end{Answer}
